\chapter{摘\texorpdfstring{\quad}{}要}
涵道式无人机凭借其结构紧凑、安全性高、气动效率优越及垂直起降能力，在低空物流、农业监测、军事侦察与应急救援等军民领域展现出广阔的应用前景。然而，涵道式无人机在近地面飞行时受到显著的地面效应扰动，导致推力与控制舵面效率随离地高度剧烈变化；同时，其复杂的气动特性与强非线性耦合使得模型参数难以离线精确标定，个体差异显著。上述问题严重制约了涵道式无人机在复杂环境下的控制精度与飞行安全。因此，开展涵道式无人机气动参数在线辨识、地面效应建模与抗扰控制研究具有重要的理论意义与工程价值。

针对气动参数难以离线获取且个体差异显著的问题，第三章提出了一种基于扩展卡尔曼滤波器的在线参数辨识方法。通过对非线性动力学模型进行合理简化与集总参数合并，构建了包含六个关键气动参数的十八维增广状态空间模型，并从理论上证明了系统的局部弱可观性。数值仿真与实际飞行试验结果表明，该算法能够在典型飞行工况下实时、稳定地估计涵道风扇基础拉力系数、侧向力系数及舵面偏转力系数等核心参数，估计不确定性区间有效覆盖真值，为后续高精度控制器设计提供了可靠的模型支撑。

针对近地面飞行时地面效应导致执行器效益非线性变化的问题，第四章开展了地面效应下的推力与控制舵面偏转力离线辨识研究。搭建了基于六维力/力矩传感器的变高度静态测试台架，系统测量了不同离地高度下涵道风扇推力与控制舵面力矩的变化规律，提出了指数衰减形式的半经验模型，并通过高斯-牛顿迭代法完成了模型参数的非线性最小二乘拟合。实验结果表明，所建模型能够准确刻画地面效应下推力的非单调变化及舵面效率随高度降低而衰减的特性，为近地面抗扰补偿提供了定量依据。

针对涵道式无人机在全飞行包线内的抗扰控制问题，第五章设计了完整的抗扰控制器。首先，基于伪逆法设计了控制分配器，实现期望力/力矩向执行器指令的映射。其次，利用实时离地高度信息动态更新控制效率矩阵中的推力系数与舵面偏转力系数，实现了地面效应的前馈补偿。静态台架与近地面实飞实验表明，补偿后推力与力矩的稳态误差显著降低，无人机在强地效区的高度与偏航角跟踪性能得到明显改善，稳态误差快速收敛且动态品质一致。进一步地，基于第三章辨识得到的气动参数，设计了非线性动态逆控制器，仿真对比表明NDI控制器能够主动解耦姿态通道间的非线性耦合，在姿态跟踪的快速性、稳定性和一致性方面优于传统串级PID控制器。

本文系统解决了涵道式无人机气动参数在线辨识、地面效应建模与抗扰控制等关键问题，所提方法在理论分析与实验验证层面均取得了积极效果，为涵道式无人机在近地面复杂环境下的高精度控制提供了有效方案。
\keywordsCN{涵道式无人机；参数辨识；地面效应；抗扰控制}

\chapter{Abstract}
Ducted-fan unmanned aerial vehicles (UAVs) exhibit broad application prospects in both military and civilian fields such as low-altitude logistics, agricultural monitoring, military reconnaissance, and emergency rescue, owing to their compact structure, high safety, superior aerodynamic efficiency, and vertical take-off and landing (VTOL) capability. However, ducted-fan UAVs suffer from significant ground effect disturbances during low-altitude flight, resulting in drastic variations in thrust and control surface efficiency with altitude above ground. Meanwhile, their complex aerodynamic characteristics and strong nonlinear coupling make it difficult to accurately calibrate model parameters offline, with considerable individual discrepancies. These issues severely restrict the control accuracy and flight safety of ducted-fan UAVs in complex environments. Therefore, research on online identification of aerodynamic parameters, ground effect modeling, and disturbance rejection control for ducted-fan UAVs is of great theoretical significance and engineering value.

To address the problem that aerodynamic parameters are hard to obtain offline and exhibit significant individual differences, Chapter 3 proposes an online parameter identification method based on the extended Kalman filter (EKF). By reasonably simplifying the nonlinear dynamic model and merging lumped parameters, an 18-dimensional augmented state-space model involving six key aerodynamic parameters is constructed, and the local weak observability of the system is theoretically proven. Numerical simulations and actual flight test results show that the algorithm can estimate core parameters such as the basic thrust coefficient, lateral force coefficient, and control surface deflection force coefficient of the ducted fan in real time and stably under typical flight conditions. The estimated uncertainty intervals effectively cover the true values, providing reliable model support for the subsequent design of high-precision controllers.

Aiming at the nonlinear variation of actuator effectiveness caused by ground effect during near-ground flight, Chapter 4 conducts offline identification research on thrust and control surface deflection forces under ground effect. A variable-altitude static test bench based on a six-axis force/torque sensor is built to systematically measure the variation laws of ducted-fan thrust and control surface moments at different altitudes above ground. A semi-empirical model in the form of exponential decay is proposed, and nonlinear least-squares fitting of model parameters is accomplished via the Gauss-Newton iteration method. Experimental results demonstrate that the established model can accurately characterize the non-monotonic variation of thrust under ground effect and the attenuation of control surface efficiency with decreasing altitude, offering a quantitative basis for near-ground disturbance rejection compensation.

For the disturbance rejection control problem of ducted-fan UAVs across the entire flight envelope, Chapter 5 designs a complete disturbance rejection controller. First, a control allocator is developed based on the pseudo-inverse method to realize the mapping from desired forces/moments to actuator commands. Second, real-time altitude data are used to dynamically update the thrust coefficients and control surface deflection force coefficients in the control effectiveness matrix, achieving feedforward compensation for ground effect. Static bench tests and near-ground flight experiments verify that the steady-state errors of thrust and moments are significantly reduced after compensation. The altitude and yaw angle tracking performance of the UAV in the strong ground-effect region is remarkably improved, with fast convergence of steady-state errors and consistent dynamic performance. Furthermore, a nonlinear dynamic inversion (NDI) controller is designed using the aerodynamic parameters identified in Chapter 3. Comparative simulation results show that the NDI controller actively decouples the nonlinear coupling between attitude channels and outperforms the conventional cascade PID controller in terms of rapidity, stability, and consistency in attitude tracking.

This work systematically solves key problems including online aerodynamic parameter identification, ground effect modeling, and disturbance rejection control for ducted-fan UAVs. The proposed methods yield positive results in both theoretical analysis and experimental verification, providing an effective solution for high-precision control of ducted-fan UAVs in complex near-ground environments.

\keywordsEN{Ducted-fan UAV; parameter identification; ground effect; disturbance rejection control}